% This is samplepaper.tex, a sample chapter demonstrating the
% LLNCS macro package for Springer Computer Science proceedings;
% Version 2.20 of 2017/10/04
%
\documentclass[runningheads]{llncs}
%
\usepackage{graphicx}
\usepackage{amsmath}
\usepackage{amssymb}
\usepackage{booktabs}
\usepackage{multirow}

\newcommand\red[1]{\color{red}#1}
% Used for displaying a sample figure. If possible, figure files should
% be included in EPS format.
%
% If you use the hyperref package, please uncomment the following line
% to display URLs in blue roman font according to Springer's eBook style:
% \renewcommand\UrlFont{\color{blue}\rmfamily}

\begin{document}

\begin{figure}
% \setlength{\belowcaptionskip}{-20pt} 
\begin{center}

\includegraphics[width=0.95\linewidth]{Fig/motivation_supp.png}
% \includegraphics[width=0.9\columnwidth]{Fig/Method.png}
\caption{
The figure depicts the experiment in the Motivation Section, which shows the SNNL of sensitive and target attributes at different inference positions for ResNet18 and VGG-11. ResNet18 and VGG-11 extract internal features from the end of each residual block and the last four max pooling layers, respectively. The $POS$ refers to the inference position in the network from shallow to deep. The results indicate a consistent decrease in the SNNL for sensitive and target attributes from shallow to deep layers across both datasets, which suggests that features are becoming more distinguishable concerning both sensitive and target attributes (accurate but unfair).
% The figure depicted the experiment in the Motivation Section, which showed the SNNL of sensitive and target attributes at different inference positions for ResNet18 and VGG-11. ResNet18 and VGG-11 extract internal features from the end of each residual block and the last four max pooling layers, respectively. The $POS$ referred to the inference position in the network from shallow to deep. The results indicated a consistent decrease in the SNNL for sensitive and target attributes from shallow to deep layers across both datasets, which suggested that features became more distinguishable concerning both sensitive and target attributes (accurate but unfair).
% We compute the SNNL of sensitive and target attributes at different inference positions and observed that the SNNL of sensitive attributes in both datasets decreased by an average of 1.1\% and 0.5\%, respectively, for each inference position from shallow to deep in ResNet18. For VGG-11, the SNNL of sensitive attributes in both datasets decreases by an average of 1.4\% and 1.2\%, respectively. Moreover, the SNNL of target attributes decreases by 0.6\% and 1.0\% for ResNet18 and by 3.7\% and 2.5\% for VGG-11. This phenomenon indicates that the features become more distinguishable with respect to both sensitive and target attributes.
}
\label{Fig.Motivation_supp}
\end{center}
\end{figure}

\begin{table*}[t]
\centering
\caption{Accuracy and fairness of classification results across different baselines with and without the ME training framework on the Fitzpatrick-17 dataset. The backbone for these experiments are VGG-11.}
\setlength{\tabcolsep}{6pt}
\label{Table:compare_fitz_me}
\scriptsize
\begin{tabular}{c c c c c c c c}\toprule
% \hline
\multicolumn{1}{c }{\text{}} & \multicolumn{1}{c }{\text{}} &\multicolumn{3}{c }{\text{Accuracy}} & \multicolumn{3}{c }{\text{Fairness}}  \\ 
\cmidrule(lr){3-5}
\cmidrule(ll){6-8}
\multicolumn{1}{ c }{\text{Method}} & \text{Skin Tone} &\text{Precision} & \text{Recall} & \text{F1-score} & \text{Eopp0 $\downarrow$} & \text{Eopp1 $\downarrow$} & \text{Eodd $\downarrow$} \\ 
\hline
\hline

\multirow{4}{*}{VGG-11} & \text{Dark} & 0.563 & 0.581 & 0.546 & \multirow{4}{*}{0.0013} & \multirow{4}{*}{0.361} & \multirow{4}{*}{0.182} \\ 
 & \text{Light} & 0.482 & 0.495 & 0.473 &  &  & \\ 
 & \text{Avg. $\uparrow$} & 0.523 & 0.538 & 0.510 &  &  & \\ 
 & \text{Diff. $\downarrow$} & 0.081 & 0.086 & 0.073 &  &  & \\ 
\hline
% \cmidrule(lr){1-1}
% \cmidrule(lr){2-2}
% \cmidrule(lr){3-5}
% \cmidrule(ll){6-8}

\multirow{4}{*}{ME-VGG-11} & \text{Dark} & 0.568 & 0.517 & 0.517 & \multirow{4}{*}{\textbf{0.0012}} & \multirow{4}{*}{\textbf{0.320}} & \multirow{4}{*}{\textbf{0.160}} \\ 
 & \text{Light} & 0.522 & 0.536 & 0.514 &  &  & \\ 
 & \text{Avg. $\uparrow$} & 0.545 & 0.527 & 0.516 &  &  & \\ 
 & \text{Diff. $\downarrow$} & 0.046 & 0.019 & 0.003 &  &  & \\ 
\hline
\hline

\multicolumn{1}{ c }{\multirow{4}{*}{HSIC}} & \text{Dark} & 0.548 & 0.522 & 0.513 & \multirow{4}{*}{0.0013} & \multirow{4}{*}{0.331} & \multirow{4}{*}{0.166} \\ 
 & \text{Light} & 0.513 & 0.506 & 0.486 &  &  & \\ 
 & \text{Avg.$\uparrow$} & 0.530 & 0.515 & 0.500 &  &  & \\ 
 & \text{Diff. $\downarrow$} & 0.040 & 0.018 & 0.029 &  &  & \\ 
\hline

\multicolumn{1}{ c }{\multirow{4}{*}{ME-HSIC}} & \text{Dark} & 0.566 & 0.528 & 0.528 & \multirow{4}{*}{\textbf{0.0012}} & \multirow{4}{*}{{\textbf{0.314}}} & \multirow{4}{*}{{\textbf{0.157}}} \\ 
 & \text{Light} & 0.523 & 0.518 & 0.503 &  &  & \\ 
 & \text{Avg. $\uparrow$} & 0.544 & 0.523 & 0.515 &  &  & \\ 
 & \text{Diff. $\downarrow$} & 0.043 & 0.010 & 0.025 &  &  & \\ 
\hline
\hline

% \multicolumn{1}{ c }{HSIC} & 0.74 & 15 & 630 & 0.70 & 15 & 618 & -\\ 
\multicolumn{1}{ c }{\multirow{4}{*}{MFD}} & \text{Dark} & 0.514 & 0.545 & 0.503 & \multirow{4}{*}{\textbf{0.0011}} & \multirow{4}{*}{0.334} & \multirow{4}{*}{0.166} \\ 
 & \text{Light} & 0.489 & 0.469 & 0.457 &  &  & \\ 
 & \text{Avg. $\uparrow$} & 0.502 & 0.507 & 0.480 &  &  & \\ 
 & \text{Diff. $\downarrow$} & 0.025 & 0.076 & 0.046 &  &  & \\ 
\hline

\multicolumn{1}{ c }{\multirow{4}{*}{ME-MFD}} & \text{Dark} & 0.561 & 0.545 & 0.535 & \multirow{4}{*}{0.0012} & \multirow{4}{*}{\textbf{0.319}} & \multirow{4}{*}{\textbf{0.159}} \\ 
 & \text{Light} & 0.539 & 0.528 & 0.512 &  &  & \\ 
 & \text{Avg. $\uparrow$} & 0.550 & 0.537 & 0.523 &  &  & \\ 
 & \text{Diff. $\downarrow$} & 0.022 & 0.017 & 0.023 &  &  & \\ 

\bottomrule
\end{tabular}
\end{table*}

\end{document}
